\chapter{CƠ SỞ LÝ THUYẾT}
\label{chap:background}

\section{Khung encoder--decoder cho bài toán sinh chuỗi}

Kiến trúc encoder--decoder tách bài toán thành hai phần: encoder nén ảnh thành
biểu diễn trung gian, decoder sinh câu mô tả một cách tự hồi quy --- từ thứ $t$
được sinh dựa trên ảnh và các từ $1..t{-}1$ đã sinh:
\begin{equation}
	p(y \mid I) = \prod_{t=1}^{T} p\!\left(y_t \mid y_{<t},\, I\right).
\end{equation}
Encoder thường là CNN huấn luyện sẵn trên ImageNet \autocite{he2016resnet};
decoder là mô hình ngôn ngữ có điều kiện. \emph{Show and Tell}
\autocite{vinyals2015show} nén cả ảnh vào một vector duy nhất --- nút thắt cổ
chai thông tin mà attention sau này giải quyết.

\section{Mạng hồi quy LSTM: giữ gradient qua chuỗi dài}

RNN thường ($h_t = \tanh(W_h h_{t-1} + W_x x_t)$) lan truyền ngược qua $T$ bước
bằng tích $T$ đạo hàm chứa thừa số $W_h^\top \operatorname{diag}(\tanh'(\cdot))$:
tích này co về 0 hoặc phát nổ theo cấp số nhân tùy trị riêng của $W_h$ --- hạn
chế cốt lõi khiến RNN khó học phụ thuộc dài. LSTM \autocite{hochreiter1997lstm}
thêm \emph{cell state} $c_t$ chạy xuyên suốt chuỗi bằng phép cộng, cùng ba cổng
sigmoid điều khiển luồng thông tin:
\begin{align}
	f_t &= \sigma(W_f [h_{t-1}; x_t] + b_f), &
	i_t &= \sigma(W_i [h_{t-1}; x_t] + b_i), \\
	o_t &= \sigma(W_o [h_{t-1}; x_t] + b_o), &
	\tilde{c}_t &= \tanh(W_c [h_{t-1}; x_t] + b_c), \\
	c_t &= f_t \odot c_{t-1} + i_t \odot \tilde{c}_t, &
	h_t &= o_t \odot \tanh(c_t),
\end{align}
với $\odot$ là nhân từng phần tử; $f_t, i_t, o_t$ là cổng quên, vào, ra. Điểm
mấu chốt: $\partial c_t/\partial c_{t-1} = f_t$ (cộng số hạng phụ) --- nhánh
cell state không có phép nhân ma trận trọng số lặp lại, chỉ nhân từng phần tử
với $f_t \in (0,1)^D$; khi cổng quên học được giá trị gần 1, gradient trên
chiều đó truyền ngược gần như không suy giảm, và các cổng cho mô hình
\emph{chọn lọc} giữ--nạp--xuất thay vì ghi đè toàn bộ trạng thái mỗi bước.

\section{Cơ chế attention}

\subsection{Attention cộng tính Bahdanau}

Thay vì một vector ảnh cố định, encoder giữ lại $R$ vùng đặc trưng
$\{f_1,\dots,f_R\}$ (ở đây $R=49$ vùng của lưới $7\times 7$). Tại mỗi bước sinh,
decoder với trạng thái ẩn $h$ tính điểm cho từng vùng bằng attention cộng tính
\autocite{bahdanau2015neural}:
\begin{equation}
	e_i = v^\top \tanh\!\left(W_h h + W_f f_i\right), \qquad
	\alpha = \operatorname{softmax}(e), \qquad
	c = \sum_{i=1}^{R} \alpha_i f_i,
	\label{eq:bahdanau}
\end{equation}
rồi dùng vector ngữ cảnh $c$ (thay đổi theo từng bước) để sinh từ tiếp theo.
\emph{Show, Attend and Tell} \autocite{xu2015show} áp dụng đúng cơ chế này cho
captioning; trọng số $\alpha$ còn cho ta bản đồ nhiệt ``từ nào nhìn vùng ảnh nào''
--- công cụ phân tích định tính chính của báo cáo.

\subsection{Self-attention và cross-attention trong Transformer}

Transformer \autocite{vaswani2017attention} thay thế hồi quy bằng attention
tích vô hướng có tỷ lệ trên toàn chuỗi:
\begin{equation}
	\operatorname{Attention}(Q,K,V)
	= \operatorname{softmax}\!\left(\frac{QK^\top}{\sqrt{d_k}}\right) V.
\end{equation}
\emph{Multi-head attention} chiếu $Q,K,V$ xuống $H$ ``đầu'' song song, mỗi đầu
một không gian con $d_k = D/H$ chiều, rồi nối kết quả và chiếu về lại $D$ chiều.
Trong đồ án, $D=512$, $H=8 \Rightarrow d_k = 64$. Trong decoder cho captioning,
mỗi khối gồm ba tầng con: (i) \emph{masked self-attention} --- chuỗi từ tự nhìn
phần quá khứ của chính nó (mặt nạ nhân quả chặn nhìn tương lai);
(ii) \emph{cross-attention} --- $Q$ đến từ chuỗi từ, $K,V$ đến từ 49 vùng ảnh,
chính là khối \emph{fusion} đa phương thức; (iii) FFN theo vị trí
($D \to d_{ff} \to D$) biến đổi phi tuyến \emph{tại từng vị trí} và chứa phần
lớn tham số của khối.

\textbf{Pre-norm và post-norm.} Bản gốc chuẩn hoá \emph{sau} mỗi sublayer
(post-norm: $x_{t+1} = \operatorname{LN}(x_t + \text{Sublayer}(x_t))$) ---
LayerNorm nằm trên đường residual chính nên gradient dễ bất ổn nếu thiếu
\emph{warmup} learning rate đủ dài. \emph{Pre-norm}
($x_{t+1} = x_t + \text{Sublayer}(\operatorname{LN}(x_t))$) giữ đường residual
thô xuyên suốt như một ``đường tắt'' gradient, ổn định hơn; đổi lại phương sai
của $x_t$ tăng theo độ sâu nên cần một LayerNorm cuối
($\text{LN}_{\text{final}}$) trước tầng đọc ra. Với dữ liệu nhỏ như Flickr8k
(30.000 caption huấn luyện), pre-norm là lựa chọn an toàn hơn.

\section{Huấn luyện teacher forcing và các kỹ thuật ổn định}

Huấn luyện chuẩn (gọi là SFT trong báo cáo này) dùng \emph{teacher forcing}:
tại bước $t$, mô hình nhận từ \emph{thật} $y_{t-1}$ làm đầu vào thay vì từ nó
vừa sinh, và cực tiểu hóa cross-entropy trung bình trên các vị trí không phải
đệm. Ngoài gradient clipping và (với Transformer) warmup + cosine decay, hai kỹ
thuật sau thay đổi trực tiếp không gian tối ưu.

\textbf{Label smoothing.} Cross-entropy với mục tiêu one-hot khuyến khích đẩy
logit của từ đúng ra vô cực --- dễ gây tự tin thái quá, đặc biệt khi dữ liệu ít
\autocite{szegedy2016rethinking}. Label smoothing trộn one-hot với phân phối
đều trên $V$ từ, hệ số $\varepsilon = 0{,}1$:
\begin{equation}
	q'(k \mid y_t) = (1-\varepsilon)\, \mathbb{1}[k = y_t] +
	\frac{\varepsilon}{V},
\end{equation}
đặt một \emph{trần} tự nhiên lên độ lớn logit, giúp gradient không bị các mẫu
``dễ'' lấn át và xác suất được hiệu chỉnh tốt hơn.

\textbf{Tie trọng số embedding--output.} Tầng embedding
($E \in \mathbb{R}^{V\times D}$) và tầng chiếu ra logits đều là ánh xạ giữa
không gian từ và không gian ẩn. \emph{Weight tying} \autocite{press2017using}
đặt $W_{\text{out}} = E^\top$: (i) giảm $D \times V$ tham số ($\approx$1,3
triệu với $V \approx 2.500$, $D = 512$); (ii) ràng buộc ngữ nghĩa --- từ có
embedding gần nhau được chiếu ra logits gần nhau, một dạng regularization hữu
ích khi dữ liệu nhỏ. Cái giá đánh đổi (thang khởi tạo của $E$ ảnh hưởng trực
tiếp đến logits) được phân tích như một lỗi thực tế gặp phải ở
mục~\ref{sec:embedding-init}.

\section{Giải mã: greedy và beam search}

Khi suy luận, mô hình tự sinh từ, dùng chính dự đoán của mình làm đầu vào bước
sau; tìm chuỗi tối ưu trên không gian $V^T$ không khả thi vét cạn nên cần giải
mã xấp xỉ. \textbf{Greedy} chọn
$\hat y_t = \arg\max_k p(k \mid \hat y_{<t}, I)$ mỗi bước, dừng khi sinh
\texttt{<eos>} hoặc chạm độ dài tối đa --- nhanh nhưng tham lam cục bộ, không
có cách quay lại sửa một lựa chọn dẫn tới ngõ cụt. \textbf{Beam search} giữ
đồng thời $B$ giả thuyết có tổng log-xác-suất cao nhất: mỗi bước mở rộng thành
$B \times V$ ứng viên rồi giữ lại $B$ ứng viên điểm cao nhất; giả thuyết sinh
\texttt{<eos>} chuyển vào tập hoàn chỉnh. Vì điểm là \emph{tổng} log-xác-suất
(mỗi số hạng $\leq 0$), so sánh trực tiếp thiên vị câu ngắn, nên điểm cuối được
chuẩn hoá theo độ dài:
\begin{equation}
	\text{score}(\hat y) = \frac{1}{|\hat y|^{\gamma}} \sum_{t=1}^{|\hat y|}
	\log p(\hat y_t \mid \hat y_{<t}, I),
\end{equation}
với $\gamma \in [0,1]$ điều chỉnh mức phạt câu dài. $B$ càng lớn không gian tìm
kiếm càng rộng ($B=1$ là greedy), nhưng chi phí tăng tuyến tính và $B$ quá lớn
trên mô hình hiệu chỉnh chưa tốt có thể chọn câu lệch khỏi phân phối ngôn ngữ
tự nhiên. Đồ án so sánh $B \in \{3, 5\}$ với greedy trên tập kiểm tra.

\section{Tinh chỉnh bằng học tăng cường: SCST}

Teacher forcing có hai điểm lệch: (i) \emph{exposure bias} --- lúc huấn luyện luôn
thấy tiền tố thật, lúc suy luận phải tự nối tiếp lỗi của chính mình; (ii) hàm mất
mát cross-entropy không trùng với thước đo đánh giá (BLEU/CIDEr).
SCST \autocite{rennie2017self} xử lý cả hai bằng REINFORCE với baseline tự phê
bình: lấy mẫu câu $y^s$, giải mã greedy câu $\hat{y}$, và cập nhật
\begin{equation}
	\nabla_\theta L \approx -\left(r(y^s) - r(\hat{y})\right)
	\nabla_\theta \log p_\theta(y^s),
\end{equation}
với $r$ là điểm CIDEr. Câu lấy mẫu tốt hơn greedy thì được khuếch đại, kém hơn
thì bị nén --- không cần huấn luyện mạng giá trị riêng.

\section{Thước đo đánh giá}

\textbf{BLEU-n} \autocite{papineni2002bleu} đo độ trùng n-gram (có hình phạt câu
ngắn) giữa câu sinh và 5 câu tham chiếu. \textbf{CIDEr} \autocite{vedantam2015cider}
đánh trọng số TF-IDF cho n-gram --- n-gram hiếm, giàu thông tin được thưởng cao
hơn --- và là thước đo tương quan tốt nhất với đánh giá của con người trong các
thước đo tự động phổ biến cho captioning.
